\documentclass[11pt,a4paper]{letter}
\usepackage[margin=1in]{geometry}
\usepackage{hyperref}
\usepackage{xcolor}

\signature{Guo Hao-Xuan \\
On behalf of the HXMT Collaboration \\
Institute of High Energy Physics, CAS \\
Beijing 100049, China \\
\href{mailto:guohx@ihep.ac.cn}{guohx@ihep.ac.cn}}

\address{Institute of High Energy Physics \\
Chinese Academy of Sciences \\
19B Yuquan Road, Shijingshan District \\
Beijing 100049, China}

\begin{document}

\begin{letter}{Editor-in-Chief \\ \emph{Science China Physics, Mechanics \& Astronomy} \\
Editorial Office, Beijing}

\opening{Dear Editor,}

We are pleased to submit our manuscript ``Saturation Recovery for Insight-HXMT/HE
Bright Burst Observations from Level-1B Raw Data'' for consideration as an
\textbf{Article} in \emph{Science China Physics, Mechanics \& Astronomy}.

The High-Energy detector (HE) on board Insight-HXMT, with its ${\sim}5100~\rm cm^2$
of effective area in the $20$--$250$~keV band, is one of the most sensitive
in-orbit instruments to bright hard X-ray transients. Its standard Level-1K
processing pipeline, however, reports gaps and constant-rate fills across
saturated intervals --- a limitation that has constrained the use of HXMT
HE data for the brightest gamma-ray bursts and for time-domain analyses of
short-timescale variability where saturation overlaps with the scientifically
most interesting moments. Until now, the most extreme HE-observed events
(GRB~200415A magnetar giant flare, GRB~221009A ultra-luminous burst, several
short and medium-bright GRBs) have been quoted in the literature with
explicit caveats about saturation-induced loss.

In this work, we present an event-level recovery pipeline that operates
\emph{directly on the Level-1B raw telemetry}, recovering both timing and
fluence in the saturation regime. The novel contributions are:

\begin{itemize}
\item A longest-increasing-subsequence (LIS) formulation for per-event
ptime ordering that is robust to the $1/16$-rate CRC ghost contamination
intrinsic to the 4-bit checksum used in HE 1B telemetry --- a structural
problem that no prior 1B reconstruction had addressed.
\item A grouped-LIS variant with a UTC-tail constraint that handles the
multi-second-event-counter wrap regime, where each event acquires multiple
candidate timestamps.
\item A cross-detector-box shape-function gap-fill, using the simultaneous
unsaturated reference-box light curve calibrated by physical area and
spectral hardness ratios, recovering not just integrated fluence but
sub-second light-curve structure.
\item An honest characterization of the method's failure modes: the
1\,ms shape-function granularity boundary (validated against ASIM/MXGS
data on GRB~200415A) and the three-box co-saturation degenerate case.
\end{itemize}

The recovered data are validated at three independent levels: against the
Level-1K pipeline on a moderately bright burst, against four external
gamma-ray instruments on the magnetar giant flare GRB~200415A, and via
within-instrument cross-box cross-correlation on GRB~221009A.
We additionally present an HE engineering-data validation that constitutes
the most stringent test available --- a controlled high-rate stream where
both saturated 1B telemetry and the unsaturated ground-truth event list
are simultaneously available.

The pipeline and associated reproduction artifacts will be released
open-source (MIT license) and archived on Zenodo at the time of
publication. During the editorial-review process the source code is
held in a private repository; access can be granted to the
editorial board and referees upon request.

\textbf{Why \emph{Science China Physics, Mechanics \& Astronomy}.} Insight-HXMT
is a flagship Chinese space astronomy mission, and a methods paper that
extends the science reach of its instruments by recovering data from one of
its most science-relevant operating regimes is strongly aligned with the
journal's scope and readership. The methodology is generalizable: the
underlying problem (FIFO-bound saturation in moderately deadtime-limited
detectors) recurs in upcoming missions, and we explicitly note in the
Conclusion how the algorithm extends to the GECAM constellation and the
SVOM/GRM payload, both of which have substantial Chinese institutional
involvement.

\textbf{Originality and submission status.} This work has not been published
elsewhere and is not under consideration by another journal. A preliminary
internal version (in Chinese) was discussed within the HXMT collaboration
but has not been distributed publicly; the present manuscript is a complete
English rewrite incorporating substantial methodological improvements and
the engineering-data validation, both unavailable at the time of the
internal version.

\textbf{Authorship.} The author list reflects the HXMT collaboration with
the corresponding-author position held by the lead developer of the
recovery pipeline. The collaboration list will be finalized in coordination
with the IHEP HXMT science operations team prior to the camera-ready
submission.

We thank you for your time considering this manuscript and look forward to
your editorial decision.

\closing{Sincerely yours,}

\end{letter}

\end{document}
